\section{The actual lattice-to-metric transport is a complete elementary chain}
\label{sec:metric-transport-complete}

\status{MT1--MT7 prove the literal implication from
\texttt{LinkDIV}, fullness, positive retained rank, \texttt{TuSize} and
\texttt{DGamSize} to \texttt{LinkMu}.  Every lattice, basis, metric,
coefficient field and finite place is the one stored in the original
\texttt{FamP}.  In particular, \texttt{NprimeDef} is a field of that
record, not an extra arithmetic hypothesis.  The actual family and the
other four inputs still have to be constructed.  This closes the separate
metric-transport leaf in the previous global audit.}

The useful simplification is that this comparison is an ideal-index
calculation followed by a rounding estimate.  It does not require a new
comparison theorem between independently constructed Hermitian spaces.
The source already keeps the two lattices inside the same metrized vector
space.  We give the entire argument, including the exact factor $2$ and
the original uniformity in the family index and word.

\subsection{The exact objects and the ideal-index identity}

% BEGIN MT1
\paragraph{MT1: the source really uses one fixed metric.}
\label{claim:MT1}
Fix a late member $(D,L)$ of an original family $\mathcal F$.  Write
$K=D.\texttt{what}_2.\texttt{LL}$, $d_K=[K:\mathbb Q]$, and let $u$
be a word.  The Hermitian-lattice record $D.\texttt{div.Nu}(u)$ consists
of a $K$-space $V_u$, a basis $b_u$, matrices $X_{u,v}$ at its infinite
places, and an $\mathcal O_K$-submodule $N_u$.  The comparison object is
\[
 L.\texttt{NuH}(u)=
       \{D.\texttt{div.Nu}(u)\ \text{with lattice }M_u\},
 \qquad M_u=L.\texttt{NuLat}(u).                         \tag{MT.1}
\]
Thus its space, basis and all infinite-place matrices are literally those
of $N_u$.  The matrices are invertible by the original
$D.\texttt{div\_good}$ field.

For the actual finite-place labels $w$, put
\[
 \ell_w=\frac{\log N\mathfrak p_w}{d_K},\qquad
 e_w=e(\mathfrak p_w/p),\qquad s_{u,w}=L.\texttt{sg}(u,w),
 \qquad q_{u,w}=\lfloor e_ws_{u,w}\rfloor.
\]
Here $\mathfrak p_w=D.\texttt{plEq}(w)$ and $p$ is its actual residue
characteristic; $D.\texttt{ram\_eq}$ identifies the stored ramification
index with $e_w$.  Define the nonzero fractional ideal
\[
       I_u=\prod_{w\in L.\texttt{supp}}\mathfrak p_w^{-q_{u,w}}.
                                                                  \tag{MT.2}
\]
The source's \texttt{NprimeDef} is exactly
\[
                N_u=I_uM_u,                                      \tag{MT.3}
\]
where the product denotes the $\mathcal O_K$-span of products.
For $x>\mathcal F.x_1$, (MT.3) is supplied by the field
$\mathcal F.\texttt{nprimeDef}(x,i)$ for every $i\in\mathcal F.I_x$.
No choice of a new lattice or an isometry is involved.
% END MT1

% BEGIN MT2
\paragraph{MT2: multiplication by an ideal changes the slope exactly.}
\label{claim:MT2}
Let $V$ be a finite-dimensional $K$-space with the fixed source metric,
$M\subset V$ a nonzero finitely generated $\mathcal O_K$-module, and
$I\ne0$ a fractional ideal.  If $r=\dim_K KM$, then
\[
       \boxed{\quad\mu(IM)=\mu(M)-\frac{2\log N(I)}{d_K}.\quad}     \tag{MT.4}
\]
This is valid for lattices in a proper subspace of $V$ as well as for
full lattices.

First let $I=\mathfrak a$ be an integral ideal.  Locally at each finite
prime, the torsion-free module $M$ is free of rank $r$ over the discrete
valuation ring.  Consequently the $\mathfrak p$-length of
$M/\mathfrak a M$ is $r\,v_{\mathfrak p}(\mathfrak a)$, and
\[
                  [M:\mathfrak a M]=N(\mathfrak a)^r.             \tag{MT.5}
\]
One may see this without a global free basis: decompose the finite
quotient into its prime-primary parts and use local bases only to count
each part.  The archimedean realization is injective, and
$\mathfrak a M$ and $M$ span the same real space, of dimension $d_Kr$.
A fundamental region for the smaller lattice consists of
$[M:\mathfrak a M]$ translates of one for the larger lattice.  Hence
its covolume is multiplied by $N(\mathfrak a)^r$.

The source normalization is
\[
 \mu(M)=-\frac{2}{d_Kr}\log\operatorname{covol}(M)
              -\log2+\frac{\log|\operatorname{disc}(K)|}{d_K}.
                                                                  \tag{MT.6}
\]
The field-dependent constant cancels, and (MT.5) gives (MT.4).
For a fractional $I$, choose $0\ne a\in\mathcal O_K$ with
$\mathfrak a=aI$ integral.  Then $a(IM)=\mathfrak a M$.
Apply the integral formula to $IM$ and $(a)$, and to $M$ and
$\mathfrak a$, and subtract.  Since
$N(I)=N(\mathfrak a)/N((a))$, this gives exactly (MT.4).

The hypotheses required in the source are slightly asymmetric: it
assumes that $N=IM$ is nonzero and finitely generated.  These imply the
same properties for $M$.  Indeed, choose $0\ne a\in I$.  Multiplication
by $a$ injects $M$ into the Noetherian module $N$, so $M$ is finitely
generated; $M=0$ would force $N=0$.  Conversely multiplication by a
nonzero fractional ideal preserves the $K$-span.  Thus the preceding
proof applies with exactly the source's stated hypotheses.
% END MT2

% BEGIN MT3
\paragraph{MT3: the entire metric discrepancy is a floor error.}
\label{claim:MT3}
Apply MT2 to (MT.3).  Because
$\log N(I_u)=-\sum_{w\in\mathrm{supp}}q_{u,w}\log N\mathfrak p_w$,
\[
 \mu(N_u)-\mu(M_u)
      =2\sum_{w\in\mathrm{supp}}\ell_w\lfloor e_ws_{u,w}\rfloor.
                                                                  \tag{MT.7}
\]
The literal local size term is
$\texttt{sz}(w,s)=2\ell_we_ws$.  It follows that
\[
 \begin{split}
 R_u&:=\mu(N_u)-\mu(M_u)
              -\sum_{w\in\mathrm{supp}}\texttt{sz}(w,s_{u,w})\\
    &=2\sum_{w\in\mathrm{supp}}\ell_w
          \bigl(\lfloor e_ws_{u,w}\rfloor-e_ws_{u,w}\bigr),
 \end{split}                                                     \tag{MT.8}
\]
and therefore
\[
      |R_u|\le2\sum_{\substack{w\in\mathrm{supp}\\s_{u,w}\ne0}}
                          \ell_w.                                \tag{MT.9}
\]
Indeed $\ell_w\ge0$, and each floor error has absolute value at most
one; it is zero if $s_{u,w}=0$.  The formula uses the ordinary integer
floor even when $s_{u,w}<0$.  The sign of $s_{u,w}$ is not a hypothesis.
In fact $R_u\le0$, although only the absolute bound is needed.
% END MT3

\subsection{The arithmetic support of the rounding error}

% BEGIN MT4
\paragraph{MT4: normalized prime-place sums and the exact support bound.}
\label{claim:MT4}
For any rational prime $p$, the degree formula
$\sum_{w\mid p}e_wf_w=d_K$ and $N\mathfrak p_w=p^{f_w}$ give
\[
                  \sum_{w\mid p}e_w\ell_w=\log p.                 \tag{MT.10}
\]
Thus for any subset of its finite places,
$\sum_{w\in S,\,w\mid p}\ell_w\le\log p$, since $e_w\ge1$.
This calculation does not assume that $K/\mathbb Q$ is Galois.  The
source's bijection $\texttt{plEq}$ and residue-characteristic identity
make (MT.10) apply to its actual labels, without duplicates.

Put
\[
 X=L.\texttt{sc}.X''(D.R.n),\quad
 T_u=(D.\texttt{mem}(u)).A.T,\quad
 \Delta_u=(D.\texttt{mem}(u)).C.\texttt{DGam},\quad
 g_0=D.\texttt{host}.g(0),\quad g_1=D.\texttt{host}.g(1),
 \quad d=D.\texttt{host}.d.
\]
The original \texttt{ZeroAbove} assertion says that above $X$, away
from $T_u,g_0,g_1,d$ and prime divisors of $\Delta_u$, both the gain
and the $E$-term vanish.  The $F$-term vanishes there by the field
$L.\texttt{F\_above}$.  Since
\[
 s_{u,w}=\sum_{j=1}^{m}
       \bigl(\texttt{gainD}(j,w)+F(j,w)-E(u,j,w)\bigr),
\]
we obtain
\[
 s_{u,w}\ne0\Longrightarrow
 p\le X\ \text{or}\ p\in\{T_u,g_0,g_1,d\}
                      \ \text{or}\ p\mid\Delta_u.                \tag{MT.11}
\]
The removal of dropped primes in \texttt{gainD} preserves the zero
gain implication.  Notice that $\Delta_u$ is not presumed prime.

For completeness, the elementary estimate
\[
       \sum_{p\le Y}\log p\le Y\log4\qquad(Y\ge0)                \tag{MT.12}
\]
can be proved by induction on $n=\lfloor Y\rfloor$.  Let
$P(n)=\prod_{p\le n}p$.  For odd $n=2a+1\ge3$, all primes in
$(a+1,2a+1]$ divide $\binom{2a+1}{a}$.  This binomial coefficient is
at most $2^{2a}=4^a$, because the two equal middle coefficients are
part of a row summing to $2^{2a+1}$.  Hence
$P(2a+1)\le P(a+1)4^a\le4^{2a+1}$.  For even $n\ge4$,
$n$ is composite, so $P(n)=P(n-1)\le4^{n-1}$; the cases $n\le2$
are immediate.  Taking logarithms proves (MT.12).  Also
$\prod_{p\mid N}p\le N$ for every positive integer $N$, so
$\sum_{p\mid N}\log p\le\log N$.

Group (MT.9) by rational residue characteristic, apply (MT.10), and
then (MT.11)--(MT.12).  Nonnegative terms allow overlaps in the union
of prime sets to be counted more than once.  This gives the full
pointwise estimate
\[
 \boxed{\quad
 |R_u|\le2\{X\log4+\log T_u+\log g_0+\log g_1+
                                      \log d+\log\Delta_u\}.\quad} \tag{MT.13}
\]
This is the source's \texttt{nprimeMu\_at}, including every exceptional
prime family and every normalization.
% END MT4

% BEGIN MT5
\paragraph{MT5: all support costs fit the original $X''$ scale.}
\label{claim:MT5}
Write $x=L.\texttt{sc}.x>1$, $\epsilon=L.\texttt{sc}.\texttt{eps}$,
$\theta=\epsilon/100$, $n=D.R.n$, and $N=D.R.N(m)$.  Recall
\[
 X'=\exp(x^{\epsilon/3}),\qquad X=X'xn,\qquad f=2n,
 \qquad m\le2n,\qquad 0<\epsilon<1/52.                            \tag{MT.14}
\]
The identity $X=X'xn$ follows directly from
$b'=a/X'$, $w''=b'/(xn)$ and $X''=a/w''$.  In particular $x\le X$.

The literal \texttt{SizeG} and \texttt{SizeQD} clauses yield
\[
 \log g_i\le x^{\theta+\epsilon}\le x\le X,
 \qquad
 \log d\le x^{\epsilon/5+\theta}\le x\le X.                       \tag{MT.15}
\]
\texttt{SizeTu} yields
\[
               \log T_u\le mx^{\theta+\epsilon}
                               \le2nx\le2X.                       \tag{MT.16}
\]
These exponent inequalities hold throughout the original permitted
$\epsilon$ range.

Here is the complete numerical argument for the remaining
$\Delta_u$ term.  The last \texttt{SizeNodes} clause gives
$Nm^{10}\le\eta/m$; since $m\ge2$, it implies $N\le\eta$.
\texttt{SizeEta} gives $\log\eta\le x^{\epsilon/4}$, so
\[
                   N\le\eta\le e^{x^{\epsilon/4}}\le X'.         \tag{MT.17}
\]
\texttt{SizeN} gives $n\le x^{2+8\epsilon+\theta}\le x^3$.
Thus $\log n\le3\log x$, while $\log N\le x^{\epsilon/4}\le x$.
Using $\log x\le x-1$ and $\log2<1$ gives
\[
                  0\le\log(fN)=\log(2nN)\le4x.
\]
Together with (MT.17) this proves
\[
                    0\le fN\log(fN)\le8X\le10X.                  \tag{MT.18}
\]
If \texttt{FactDGam} holds with constant $C_\Delta$, it says exactly
$\log\Delta_u\le C_\Delta fN\log(fN)$.  Hence
\[
                   \log\Delta_u\le10|C_\Delta|X.                 \tag{MT.19}
\]
The absolute value permits the constant in the original record to have
either sign.  Inserting (MT.15)--(MT.19) into (MT.13) proves
\[
                  |R_u|\le
          2\bigl(\log4+5+10|C_\Delta|\bigr)X.                    \tag{MT.20}
\]
More generally, if $\log T_u\le C_TX$ and
$\log\Delta_u\le C_\Gamma X$, the exact constant is
$2(\log4+3+C_T+C_\Gamma)$, as in \texttt{nprimeMuR\_at}.
% END MT5

\subsection{Uniform installation on the original family}

% BEGIN MT6
\paragraph{MT6: literal \texttt{LinkMu} from the original source inputs.}
\label{claim:MT6}
Let $\mathcal F:\texttt{FamP}(\epsilon)$ with $\epsilon>0$.  Suppose
\[
 \texttt{LinkDIV}(\mathcal F,C_{\rm abs},C_H),\qquad
 \texttt{NuFullF}(\mathcal F),\qquad
 \texttt{J.MuRank}(\mathcal F),\qquad
 \texttt{TuSize}(\mathcal F),\qquad
 \texttt{DGamSize}(\mathcal F).                                  \tag{MT.21}
\]
Then the literal \texttt{LinkMu} holds: there are one constant $C$
and one threshold $x_0$ such that for every real $x\ge x_0$, every
$i\in\mathcal F.I_x$ and every word $u$,
\[
 \left|\mu(N_u)-\mu(M_u)
       -\sum_{w\in\mathrm{supp}}\texttt{sz}(w,s_{u,w})\right|
                 \le C X''(D.R.n).                              \tag{MT.22}
\]

To prove this with precisely the recorded quantifiers, first extract the
threshold from each input in (MT.21).  \texttt{LinkDIV} supplies
\texttt{ZeroAbove} for all late $x,i$.  \texttt{NuFullF} supplies finite
generation of every $N_u$; \texttt{MuRank} gives its positive $K$-rank,
hence $N_u\ne0$.  The metric matrices remain invertible by
$D.\texttt{div\_good}$, and (MT.3) is already a field of
$\mathcal F$ once $x>\mathcal F.x_1$.

\texttt{TuSize} initially bounds
$\log T_u\le\mathcal F.C_{Tu}\,m x^\epsilon$, with a family
constant $\mathcal F.C_{Tu}$.  Since $x^{\epsilon/100}\to\infty$,
one further threshold makes
$\mathcal F.C_{Tu}\le x^{\epsilon/100}$.  This gives the literal
\texttt{SizeTu} inequality, and hence (MT.16), uniformly in $i,u$.
\texttt{DGamSize} supplies one $C_\Delta$ before $x,i,u$ and the
literal \texttt{FactDGam} bound.  Every size assertion used in MT5
belongs to $\mathcal F.\texttt{sizes}(x,i)$.

Take the maximum of these finitely many thresholds and
$\mathcal F.x_1+1$.  MT1--MT5 now apply simultaneously to every $i,u$.
Take $C=2(\log4+5+10|C_\Delta|)$.  This proves (MT.22), with no
finiteness assumption on the union of all index sets or on the family
of words as $x$ varies.  The thresholds were already uniform in those
variables before their maximum was taken.

If one starts instead with the source's \texttt{MuTu} and
\texttt{MuDGam}, the same proof uses their supplied constants $C_T$
and $C_\Gamma$ directly and gives
$C=2(\log4+3+C_T+C_\Gamma)$.  This is exactly
\texttt{J.linkMu\_of\_linkDIV}; the preceding two paragraphs also
supply its entire \texttt{TuSize}/\texttt{DGamSize} installation.
% END MT6

% BEGIN MT7
\paragraph{MT7: a smaller complete global interface.}
\label{claim:MT7}
In the source's \texttt{topDesWin} theorem, \texttt{DesignExists}
chooses a single $\mathcal F,C_H$ after the finite threshold join.
It supplies \texttt{ChosenE}, \texttt{TuSize} and \texttt{DGamSize}.
The closed rank/fullness input gives
\texttt{NuFullF}$\wedge$\texttt{MuRank} on this same $\mathcal F$.
As in DA2, \texttt{HypV}, \texttt{LinkN} and
\texttt{SmallLevelRankEq} give \texttt{DIVinN}; then
\texttt{LinkEta} and \texttt{LinkDIVt} give \texttt{LinkDIV}.
Apply MT6.  The resulting \texttt{LinkMu} is on exactly the family
already carrying the other four \texttt{GlobalTheoremInputs}.
Consequently DA2's metric-transport premise is fully discharged.

The sufficient global interface is therefore the original design
existence, the nine original closed link/rank conditions, and the
original class inclusions and \texttt{HypV}.  No independent
\texttt{METRIC\_TRANSPORT} leaf is needed.  This is a genuine reduction
in the human-proof frontier, although it does not construct
\texttt{DesignExists}, \texttt{LinkDIVt}, or the rank/fullness input.
In particular, actual spectral/Gram bounds in \texttt{LinkS} and
height/operator bounds in \texttt{LinkDIV} keep their independent roles.
% END MT7
